\documentclass[12pt]{article}
\usepackage[T1]{fontenc}
\usepackage[utf8]{inputenc}
\usepackage{lmodern}
\usepackage{textcomp}
\usepackage{enumitem}
\usepackage{geometry}
\usepackage{graphicx}
\usepackage{amsmath, amssymb}
\geometry{margin=1in}
\begin{document}
\begin{center}
Constitutional Law - Graduate Level Assessment 19 \
Total Marks: 40 \
Answer all questions.
\end{center}

\section*{Multiple Choice (10 marks)}
\begin{enumerate}[label=\arabic*.]
\item Which one of the following statements is false?
\begin{enumerate}[label=(\alph*)]
    \item Protocol 14 established the Human Rights Council
    \item Protocol 14 mandated a two judge formation for hearing admissibility cases
    \item Protocol 14 reduced the jurisdiction of the European Court of Human Rights
    \item Protocol 14 shortened the judicial term of office for the European Court of Human Rights
    \item Protocol 14 introduced a requirement for unanimous decision-making in admissibility cases
\end{enumerate}
\item A defendant charged with first-degree murder shall be furnished with a list containing names and addresses of all prospective jurors
\begin{enumerate}[label=(\alph*)]
    \item only if the prosecutor agrees.
    \item immediately after being charged.
    \item upon request and showing of good cause.
    \item upon court order.
    \item under no circumstances.
\end{enumerate}
\item The most likely contract to be classified under the Uniform Commercial Code (UCC) is a contract for
\begin{enumerate}[label=(\alph*)]
    \item The purchase of a commercial property building.
    \item A contract for hiring a live-in nanny.
    \item The purchase of stocks and bonds.
    \item Crops and timber to be severed from the property next summer.
    \item A common carrier delivering a new computer.
\end{enumerate}
\item The \_\_\_\_\_\_\_\_ School of jurisprudence believes that the law is an aggregate of social traditions and customs that have developed over the centuries.
\begin{enumerate}[label=(\alph*)]
    \item Command
    \item Sociological
    \item Historical
    \item Interpretive
    \item Analytical
\end{enumerate}
\item Which of the following quotations best describes the central thesis of difference feminism?
\begin{enumerate}[label=(\alph*)]
    \item Men are unable to comprehend their differences from women.'
    \item Men and women have different conceptions of the feminist project.'
    \item Women look to context, whereas men appeal to neutral, abstract notions of justice.'
    \item 'Women and men have identical perspectives on justice.'
    \item 'Feminism is about eradicating differences between men and women.'
\end{enumerate}
\end{enumerate}

\section*{True / False (10 marks)}
\textit{Answer True or False}
\begin{enumerate}[label=\arabic*.]
\setcounter{enumi}{5}
\item True or False: The Analytical School posits that a custom becomes law upon receiving judicial recognition from the courts.
\item True or False: Hart's distinction primarily differentiates between legal and ethical obligations within the context of law and morality.
\item True or False: The use of armed force prior to the United Nations Charter was primarily permitted only in cases of self-defense.
\item True or False: The UK Constitution is uncodified and is derived from various sources including statutes, conventions, and judicial decisions.
\item True or False: Standard-setting in human rights diplomacy exclusively involves proposing binding legal standards that must be adhered to by states.
\end{enumerate}

\section*{Long Form Response (20 marks)}
\textit{Answer each question with a clear and complete explanation.}
\begin{enumerate}[label=\arabic*.]
\setcounter{enumi}{10}
\item Discuss the implications of the Kadi judgment on the incorporation of UN Security Council resolutions in national law, particularly in relation to human rights protections.
\item Discuss the factors that may constitute sufficient good cause or excusable neglect for a litigant to file a late notice of appeal, using a specific scenario to illustrate your analysis.
\end{enumerate}

\vfill
\noindent\textit{End of Paper}
\end{document}